\section{Benchmark Output Parsing and Output Structure}

The benchmark execution pipeline produces traceable and independently analyzable artifacts. All generated data is stored under \texttt{Benchmark\_Framework/outputs}, which is divided according to the main stages of the evaluation: generation, parsing, and scoring.

This separation is a core design choice. It makes it possible to distinguish between different failure modes: a model may generate text that contains a reasonable plan but is not cleanly parseable, or it may produce a syntactically valid action sequence that still fails during formal validation. The artifacts are persisted when a benchmark case terminates, either because a valid plan has been found or because the configured maximum number of iterations has been reached.

The \texttt{outputs} directory is organized into three main subdirectories:
\begin{itemize}
    \item \textbf{raw}: stores the direct model interaction and generation metadata.
    \item \textbf{parsed}: stores the structured action sequences extracted from the raw text.
    \item \textbf{scored}: stores validation outcomes, repair information, and aggregate metrics for the case.
\end{itemize}

Each layer follows the same hierarchical path:
\begin{center}
\texttt{<run\_id>/<model>/<protocol>/<domain>/<difficulty>/<instance>.json}
\end{center}
This layout makes the three artifacts for the same benchmark instance directly comparable across the raw, parsed, and scored layers.

\subsection{Raw}

The files in \texttt{outputs/raw/} are the closest record of the interaction with the LLM. Each artifact stores the prompts sent to the model, the model response, any provider-side reasoning text when available, and generation metadata such as token usage, latency, and notes.

This layer is the baseline record of what the model produced before any interpretation by the parser or validator.

\subsubsection{Raw File Structure}

Listing~\ref{lst:raw_json_structure} shows the main fields stored in a raw artifact. Provider-specific generation fields may vary across backends, but the surrounding structure is shared.

\begin{lstlisting}[language=Python, caption={Schematic structure of a raw JSON artifact}, label={lst:raw_json_structure}]
{
  "model_id": "hf_gemma_4_31b_it",
  "task_id": "block-grouping_easy_pfile1",
  "protocol_id": "iterative_repair",
  "task_family": "block-grouping",
  "tier": "easy",
  "instance_id": "pfile1",
  "generation_time_seconds": 3526.41,
  "attempts": [
    {
      "iteration": 1,
      "generation_time_seconds": 623.25,
      "messages": [
        {
          "role": "system",
          "content": "You are a PDDL planning assistant..."
        },
        {
          "role": "user",
          "content": "TASK FAMILY: block-grouping..."
        }
      ],
      "generation": {
        "raw_text": "(move_block_left b1)\n(move_block_left b1)\n...",
        "reasoning_text": "Reasoning text, when available...",
        "usage": {
          "prompt_tokens": 2125,
          "completion_tokens": 4096,
          "total_tokens": 6221
        },
        "latency_s": 536.60,
        "notes": [],
        "generation_time_seconds": 623.25
      }
    }
  ]
}
\end{lstlisting}

\subsection{Parsed}

The \texttt{outputs/parsed/} directory contains the result of the first processing step after generation. The parser attempts to extract executable PDDL actions from the raw model text and stores the extracted plan in a structured form.

Each parsed artifact includes:
\begin{itemize}
    \item the sequence of actions extracted from the visible model answer;
    \item parser diagnostics such as missing actions, unknown action names, wrong arity, markdown removal, or compressed-action expansion;
    \item a separate reasoning-derived action sequence when provider-side reasoning text contains a candidate plan;
    \item a pointer to the corresponding raw artifact.
\end{itemize}

This layer bridges the gap between unstructured natural language and the action sequence that can be submitted to validation. If no usable action can be extracted, the benchmark records a parser failure and the case can still be scored without invoking VAL on an empty plan.

\subsubsection{Parsed File Structure}

Listing~\ref{lst:parsed_json_structure} shows the current parsed artifact schema. The \texttt{raw} subsection represents the official plan extracted from the visible answer; the \texttt{reasoning} subsection is diagnostic and is used to study whether a plan was present only in reasoning text.

\begin{lstlisting}[language=Python, caption={Schematic structure of a parsed JSON artifact}, label={lst:parsed_json_structure}]
{
  "model_id": "hf_gemma_4_31b_it",
  "task_id": "block-grouping_easy_pfile1",
  "protocol_id": "iterative_repair",
  "task_family": "block-grouping",
  "tier": "easy",
  "instance_id": "pfile1",
  "generation_time_seconds": 3526.41,
  "raw_output_path": ".../outputs/raw/<run_id>/<model>/<protocol>/<domain>/<difficulty>/pfile1.json",
  "attempts": [
    {
      "iteration": 1,
      "generation_time_seconds": 623.25,
      "parsed_plan": {
        "raw": {
          "actions": [
            "(move_block_left b1)",
            "(move_block_left b1)"
          ],
          "format_issues": [
            "raw_text_contains_reasoning_like_content"
          ],
          "contains_reasoning": true,
          "source_kind": "reasoning_like_raw"
        },
        "reasoning": {
          "actions": [
            "(move_block_left b1)",
            "(move_block_left b1)"
          ],
          "format_issues": [
            "multiple_reasoning_candidate_plans"
          ],
          "source_ref": {
            "artifact": "raw",
            "field": "generation.reasoning_text"
          }
        }
      }
    }
  ]
}
\end{lstlisting}

\subsection{Scored}

The final layer, \texttt{outputs/scored/}, contains the normalized outcome of the case. If the parser extracts at least one action, the action sequence is validated with VAL. If parsing fails, the scored artifact stores a normalized \texttt{parse\_error} result instead. In both cases, the scored layer provides the fields used by the analysis notebooks and report metrics.

Each scored artifact includes:
\begin{itemize}
    \item the final \texttt{solved} decision for the benchmark instance;
    \item the number of iterations used and whether the iteration limit was reached;
    \item aggregate metrics such as \texttt{validity\_at\_1}, \texttt{validity\_at\_k}, plan length, repair success, and failure type;
    \item one validation-style record for each attempt;
    \item repair feedback and reasoning-plan diagnostics when the iterative protocol or reasoning extraction uses them.
\end{itemize}

\subsubsection{Scored File Structure}

Listing~\ref{lst:scored_json_structure} shows the main scored fields. The VAL-specific validation fields are nested under \texttt{validation\_result}; prefix-validation, reasoning, and feedback fields belong to the attempt record itself.

\begin{lstlisting}[language=Python, caption={Schematic structure of a scored JSON artifact}, label={lst:scored_json_structure}]
{
  "model_id": "hf_gemma_4_31b_it",
  "task_id": "block-grouping_easy_pfile1",
  "protocol_id": "iterative_repair",
  "task_family": "block-grouping",
  "tier": "easy",
  "instance_id": "pfile1",
  "solved": false,
  "iterations_used": 8,
  "max_iterations": 8,
  "stopped_by_iteration_limit": true,
  "generation_time_seconds": 3526.41,
  "metrics": {
    "validity_at_1": false,
    "validity_at_k": false,
    "repair_success": false,
    "iterations_to_valid": null,
    "plan_length": 30,
    "error_type": "unsatisfied_goal",
    "hit_iteration_limit": true
  },
  "raw_output_path": ".../outputs/raw/<run_id>/<model>/<protocol>/<domain>/<difficulty>/pfile1.json",
  "parsed_output_path": ".../outputs/parsed/<run_id>/<model>/<protocol>/<domain>/<difficulty>/pfile1.json",
  "attempts": [
    {
      "iteration": 1,
      "generation_time_seconds": 623.25,
      "validation_result": {
        "valid": false,
        "status": "invalid",
        "error_type": "unsatisfied_goal",
        "feedback_text": "The plan finishes without achieving the goal...",
        "failed_step": null,
        "failed_action": null,
        "goal_satisfied": false,
        "plan_length": 30,
        "validation_time_ms": 11,
        "raw_validator_output": "Checking plan...",
        "details": {
          "validator_kind": "external",
          "return_code": 1
        }
      },
      "first_valid_prefix_length": null,
      "first_valid_plan_text": null,
      "final_plan_valid": false,
      "extra_actions_after_first_valid": null,
      "reasoning_validation_result": null,
      "reasoning_first_valid_prefix_length": null,
      "reasoning_first_valid_plan_text": null,
      "reasoning_final_plan_valid": null,
      "reasoning_extra_actions_after_first_valid": null,
      "reasoning_candidate_count": null,
      "reasoning_valid_candidate_count": null,
      "reasoning_selected_candidate_index": null,
      "reasoning_selected_candidate_truncated": null,
      "feedback_to_next_iteration": "Use the validator feedback as a diagnostic hint..."
    }
  ]
}
\end{lstlisting}
